\documentclass[10pt,a4paper]{article}
\usepackage[utf8]{inputenc}
\usepackage[T1]{fontenc}
\usepackage{lmodern}
\usepackage[spanish,es-nodecimaldot]{babel}
\usepackage[a4paper,margin=2.1cm,headheight=14pt]{geometry}
\usepackage{booktabs,longtable,array,amsmath,enumitem,xcolor,hyperref,tikz}
\usetikzlibrary{arrows.meta,positioning}
\usepackage{fancyhdr}
\usepackage{microtype}

\hypersetup{
  colorlinks=true,
  linkcolor=blue!45!black,
  urlcolor=blue!55!black,
  pdftitle={Especificación de Inteligencia de Negocios - AdventureWorks},
  pdfauthor={AdventureWorks}
}
\pagestyle{fancy}
\fancyhf{}
\lhead{AdventureWorks | Inteligencia de Negocios}
\rhead{Decisiones D2 y D7}
\cfoot{\thepage}
\setlist{nosep,leftmargin=1.5em}
\renewcommand{\arraystretch}{1.12}
\setlength{\tabcolsep}{3pt}
\setlength{\parskip}{0.35em}
\setlength{\parindent}{0pt}
\setlength{\emergencystretch}{2em}
\sloppy

\definecolor{navy}{RGB}{25,54,93}
\definecolor{pale}{RGB}{231,239,248}
\definecolor{barone}{RGB}{72,126,180}
\definecolor{bartwo}{RGB}{104,160,119}
\definecolor{barthree}{RGB}{227,166,73}
\tikzset{
  box/.style={draw=navy,rounded corners,fill=pale,align=center,
    text width=2.45cm,minimum height=0.8cm,font=\small},
  reqboxa/.style={draw=navy,rounded corners,fill=pale,align=center,
    text width=2.7cm,minimum height=1.35cm,font=\small},
  reqboxb/.style={draw=navy,rounded corners,fill=pale,align=center,
    text width=3.0cm,minimum height=1.35cm,font=\small},
  reqboxc/.style={draw=navy,rounded corners,fill=pale,align=center,
    text width=4.2cm,minimum height=1.35cm,font=\small},
  flow/.style={-{Stealth[length=2mm]},thick,draw=navy}
}

\begin{document}
\hypersetup{pageanchor=false}

\begin{titlepage}
\centering
\vspace*{2cm}
{\Large\textsc{AdventureWorks Cycles}\par}
\vspace{1.2cm}
{\huge\bfseries Especificación de Inteligencia de Negocios\par}
\vspace{0.6cm}
{\LARGE Decisiones D2 y D7\par}
\vspace{1.2cm}
{\large Requerimientos informacionales y diseño del subconjunto del Data Warehouse\par}
\vfill
\end{titlepage}
\hypersetup{pageanchor=true}

\section{Objetivo y alcance}

Esta especificación define la información necesaria para analizar dos
decisiones de AdventureWorks: lanzar una oferta especial para un producto
(D2) e incrementar las existencias de una materia prima (D7). Ambas se
relacionan con el objetivo de aumentar los márgenes de ganancia mediante el
crecimiento de ingresos y la gestión de costos.

El documento organiza los requerimientos del negocio, las fuentes y sus
modelos, la arquitectura informacional, las visualizaciones y las reglas de
transformación y carga. Las definiciones se apoyan en las tablas operacionales
de AdventureWorks y los ejemplos utilizan registros de sus archivos de datos.
La implementación aplica estas estructuras y reglas a la versión de las
fuentes configurada para el entorno.

\section{Requerimientos de negocio}

\textbf{Resumen.} Las decisiones D2 y D7 conectan el objetivo de aumentar los
márgenes de ganancia con el crecimiento de ingresos y la gestión de costos.
Los indicadores y dimensiones propuestos describen los datos que permiten
analizar cada decisión.

\textbf{Relación entre objetivo, decisiones e información.}
\begin{center}
\begin{tikzpicture}[node distance=0.65cm]
  \node[reqboxa] (goal) {Objetivo\\Aumentar los márgenes de ganancia};
  \node[reqboxb,right=of goal] (decisions)
    {Decisiones\\D2: oferta especial\\D7: existencias de materia prima};
  \node[reqboxc,right=of decisions] (information)
    {Información requerida\\Ventas, descuentos, costos y recompra\\
     Inventario, umbrales, compras y proveedores};
  \draw[flow] (goal) -- (decisions);
  \draw[flow] (decisions) -- (information);
\end{tikzpicture}
\end{center}

\subsection*{D2 — Lanzar una oferta especial para un producto}

\textbf{Descripción y propósito.} Determinar si el descuento temporal en un
producto produce suficiente volumen e ingreso neto para justificar la
reducción del precio. La oferta puede también atraer compradores que regresen
posteriormente a precio regular. La evaluación integra unidades e ingreso neto para contrastar la respuesta
comercial con el descuento aplicado.

\textbf{Justificación de selección.} La decisión conecta directamente la
política promocional con unidades, ingreso neto, margen bruto y recompra.

\textbf{Preguntas de negocio}
\begin{enumerate}
  \item ¿Qué productos y ofertas generan más unidades e ingreso neto durante
  su vigencia y qué margen bruto estimado obtienen?
  \item ¿Qué descuento efectivo se aplicó y cómo se comparan las ventas por
  día con un período base comparable?
  \item ¿Cómo varían los resultados por territorio, canal de venta y tipo de
  cliente?
  \item ¿Qué proporción de clientes que compró con descuento realizó después
  otra compra a precio regular?
\end{enumerate}

\textbf{Indicadores.} En las fórmulas, \(Q\) es cantidad, \(P\) precio unitario,
\(d\) descuento efectivo, \(L\) total neto de línea, \(C_{std}\) costo estándar
y \(D\) número de días.
\texttt{LineTotal} es el neto transaccional. El margen bruto estimado utiliza
el costo estándar válido en la fecha del pedido y compara importes convertidos
a USD mediante la tasa asociada al pedido.

\small
\begin{longtable}{p{0.15\textwidth}p{0.24\textwidth}p{0.23\textwidth}p{0.11\textwidth}p{0.16\textwidth}}
\toprule
\textbf{KPI} & \textbf{Definición y propósito} & \textbf{Fórmula} &
\textbf{Unidad} & \textbf{Fuente / pregunta}\\
\midrule
\endhead
Unidades con oferta & Volumen vendido en líneas con descuento aplicado;
dimensiona la respuesta observada. & \(\sum Q\), \(d>0\) & unidades &
SalesOrderDetail (SOD); P1, P2\\
Venta neta con oferta & Ingreso neto asociado a esas líneas. & \(\sum L\),
por moneda & moneda & SOD; P1\\
Margen bruto estimado & Contrasta ingreso neto y costo estándar vigente para
analizar el aporte comercial. & \(\frac{\sum (L_{USD}-Q\cdot C_{std})}
{\sum L_{USD}}\), líneas con costo y tasa asociados & porcentaje &
SOD, ProductCostHistory, CurrencyRate; P1, P2\\
Descuento efectivo ponderado & Reducción media del precio antes del
descuento, ponderada por unidades. & \(\frac{\sum QPd}{\sum QP}\) &
porcentaje & SOD; P2\\
Variación de unidades por día & Compara el volumen de oferta y base ajustando
por duración. & \(\frac{Q_{oferta}/D_{oferta}}
{Q_{base}/D_{base}}-1\) & porcentaje & SOD, SalesOrderHeader (SOH); P2\\
Recompra regular observada & Clientes que vuelven después de una línea con
descuento y compran sin descuento. & \(\frac{clientes\ con\ recompra}
{clientes\ con\ oferta}\) & porcentaje & SOH, SOD, Customer; P4\\
\bottomrule
\end{longtable}
\normalsize

La comparación de períodos utiliza el mismo producto, moneda, canal y duración
para facilitar la lectura de las ventas observadas antes, durante y después
de la oferta.

\textbf{Parámetros.} Fecha y ventana base para delimitar la vigencia y comparar
tasas diarias; producto/categoría y oferta para seleccionar la decisión;
territorio y canal para identificar diferencias de mercado; tipo de cliente
(persona o revendedor) para separar segmentos; moneda para evitar agregaciones
incompatibles. Cada filtro permite acotar las comparaciones de P1--P4.

\subsection*{D7 — Incrementar la cantidad almacenada de una materia prima}

\textbf{Descripción y propósito.} Evaluar el incremento de inventario de un
material como medida de cobertura ante cambios de precio o suministro. El
análisis integra las existencias, los umbrales de planificación y el historial
de precios y recepciones de compra. El punto de reorden y el stock de seguridad
de Product se aplican al total de existencias del producto; la distribución
por ubicación aporta contexto operativo.

\textbf{Justificación de selección.} La decisión relaciona capital de
inventario, necesidades de abastecimiento y condiciones de compra, vinculando
las existencias con el objetivo de gestionar costos.

\textbf{Preguntas de negocio}
\begin{enumerate}
  \item ¿Qué materiales están por debajo de su punto de reorden o stock de
  seguridad, considerando su existencia total?
  \item ¿Cómo se distribuye esa existencia entre las ubicaciones registradas?
  \item ¿Cómo han variado los precios de compra recientes y qué proveedores
  concentran las cantidades recibidas?
  \item ¿Qué cantidades fueron ordenadas, recibidas y rechazadas para los
  materiales priorizados?
\end{enumerate}

\textbf{Indicadores.} \(H\) representa la existencia total del producto,
\(R_p\) su punto de reorden, \(S_p\) su stock de seguridad, \(q_r\) cantidad
recibida y \(c\) proveedor.
POD = PurchaseOrderDetail; POH = PurchaseOrderHeader.

\small
\begin{longtable}{p{0.15\textwidth}p{0.24\textwidth}p{0.23\textwidth}p{0.11\textwidth}p{0.16\textwidth}}
\toprule
\textbf{KPI} & \textbf{Definición y propósito} & \textbf{Fórmula} &
\textbf{Unidad} & \textbf{Fuente / pregunta}\\
\midrule
\endhead
Existencia total & Cantidad registrada para el producto sumando bins y
ubicaciones. & \(H=\sum Quantity\) & unidades & ProductInventory (PI);
P1, P2\\
Brecha al umbral & Cantidad que falta para alcanzar el nivel seleccionado;
prioriza revisión. & \(\max(0,R_p-H)\) o \(\max(0,S_p-H)\) & unidades &
Product; P1\\
Precio recibido ponderado & Precio medio de compra ponderado por unidades
recibidas en la ventana. & \(\frac{\sum UnitPrice\cdot q_r}{\sum q_r}\) &
moneda por unidad & POD; P3\\
Variación del precio recibido & Compara el precio ponderado entre meses
consecutivos con recepciones. & \(\frac{Precio_{reciente}}{Precio_{base}}-1\) &
porcentaje & POH; POD; P3\\
Porcentaje rechazado & Señala la fracción de cantidad recibida que fue
rechazada. & \(\frac{\sum RejectedQty}{\sum ReceivedQty}\) & porcentaje &
POD; P4\\
Participación del proveedor principal & Mide concentración de unidades
recibidas por material en la ventana. & \(\frac{\max_c(\sum q_{r,c})}
{\sum_c q_{r,c}}\) & porcentaje & POH, POD, Vendor; P3, P4\\
\bottomrule
\end{longtable}
\normalsize

La brecha frente al umbral orienta la priorización del análisis de inventario
junto con precios, recepciones y participación de proveedores.

\textbf{Parámetros.} Producto/material para acotar la decisión; fecha/ventana
de compra para medir precios y recepciones; ubicación para analizar la
distribución física; proveedor para evaluar concentración; umbral elegido
(reorden o seguridad) para calcular la brecha.

\subsection*{Visualizaciones propuestas}

\small
\begin{longtable}{p{0.12\textwidth}p{0.23\textwidth}p{0.22\textwidth}p{0.35\textwidth}}
\toprule
\textbf{Vista} & \textbf{Pregunta y KPI} & \textbf{Visualización} &
\textbf{Dimensiones, filtros y justificación}\\
\midrule
\endhead
D2.1 & ¿Cómo cambian unidades, venta neta y margen durante la oferta? Unidades,
venta neta, margen estimado y descuento. & Línea temporal por día o semana. &
Fecha, producto y oferta; comparar con ventana base de duración equivalente.
Facilita contrastar ventas antes, durante y después.\\
D2.2 & ¿Qué oferta/producto rinde mejor? Unidades, venta neta, margen y
descuento efectivo. & Tabla ordenable con indicadores destacados. & Producto,
categoría, oferta, moneda y período. Presenta ingresos, margen y descuento
en indicadores separados.\\
D7.1 & ¿Qué materiales requieren revisión? Existencia total y brecha al
umbral. & Tabla de excepciones con barras de estado. & Producto y umbral; las
barras muestran existencias por ubicación y el indicador de brecha se calcula
contra el total del producto. Facilita priorizar materiales.\\
D7.2 & ¿Qué proveedor/precio acompaña el material? Precio ponderado,
variación y participación de proveedor. & Línea de precio por período y
proveedor, con tabla de cantidades recibidas/rechazadas. & Producto,
proveedor y ventana. La tendencia y el volumen juntos dan contexto al precio.\\
\bottomrule
\end{longtable}
\normalsize

\textbf{Ejemplos con filas de AdventureWorks.} Los siguientes registros
ilustran el tipo de análisis que habilitan las vistas propuestas.

\begin{itemize}
  \item \textbf{D2: catálogo de oferta.} \texttt{SpecialOffer.csv} contiene
  \emph{Mountain-100 Clearance Sale} (ID 7, 35\%), \emph{Road-650 Overstock}
  (ID 9, 30\%) y \emph{Mountain Tire Sale} (ID 10, 50\%). Estos porcentajes
  son descuentos nominales; el descuento realizado se mide en SOD.
  \item \textbf{D2: línea base.} En SOD, la orden 43659, línea 1, registra el
  producto 776, una unidad, precio 2024.9940, descuento 0.0000 y total neto
  2024.9940. Este registro ejemplifica una venta regular para comparar con
  períodos de oferta del mismo producto y moneda.
  \item \textbf{D7: inventario por ubicación.} Product 317 (\emph{LL Crankarm})
  tiene 283 unidades en Location 1, 158 en Location 5 y 152 en Location 50:
  total 593. Su ReorderPoint es 375 y SafetyStockLevel es 500. La comparación
  agregada sitúa la existencia por encima de ambos umbrales, con el detalle
  por ubicación disponible para analizar la distribución.
  \item \textbf{D7: compra.} La línea 8 de la orden de compra 7 corresponde al
  producto 317: 550 unidades ordenadas y recibidas, ninguna rechazada, precio
  unitario 27.0585 y total de línea 14882.1750. La fecha prevista registrada
  es 13 de mayo de 2022.
\end{itemize}

\begin{figure}[htbp]
\centering
\begin{tikzpicture}[x=0.10cm,y=0.72cm]
  \draw[-{Stealth},navy] (0,0) -- (65,0) node[right,font=\small]{descuento nominal (\%)};
  \foreach \x in {0,10,...,60} {
    \draw[gray!45] (\x,0) -- (\x,-0.12);
    \node[below,font=\scriptsize] at (\x,-0.12) {\x};
  }
  \node[anchor=east,font=\small] at (-1,2.1) {Mountain-100 Clearance Sale};
  \node[anchor=east,font=\small] at (-1,1.1) {Road-650 Overstock};
  \node[anchor=east,font=\small] at (-1,0.1) {Mountain Tire Sale};
  \fill[barone] (0,1.85) rectangle (35,2.35);
  \fill[bartwo] (0,0.85) rectangle (30,1.35);
  \fill[barthree] (0,-0.15) rectangle (50,0.35);
  \node[anchor=west,font=\small] at (36,2.1) {35\%};
  \node[anchor=west,font=\small] at (31,1.1) {30\%};
  \node[anchor=west,font=\small] at (51,0.1) {50\%};
\end{tikzpicture}
\caption{Descuentos nominales registrados para tres ofertas en SpecialOffer.}
\end{figure}

\begin{figure}[htbp]
\centering
\begin{tikzpicture}[x=0.010cm,y=0.8cm]
  \draw[-{Stealth},navy] (0,0) -- (690,0) node[right,font=\small]{unidades};
  \foreach \x in {0,100,...,600} {
    \draw[gray!45] (\x,0) -- (\x,-0.12);
    \node[below,font=\scriptsize] at (\x,-0.12) {\x};
  }
  \fill[barone] (0,0.18) rectangle (283,0.68);
  \fill[bartwo] (283,0.18) rectangle (441,0.68);
  \fill[barthree] (441,0.18) rectangle (593,0.68);
  \draw[very thick,red!70!black] (375,0.02) -- (375,1.15);
  \draw[very thick,orange!75!black,dashed] (500,0.02) -- (500,1.85);
  \node[font=\small,anchor=south] at (375,1.15) {Reorden: 375};
  \node[font=\scriptsize,anchor=south] at (500,1.95) {Seguridad: 500};
  \node[font=\scriptsize,white] at (141.5,0.43) {Loc. 1: 283};
  \node[font=\scriptsize,white] at (362,0.43) {Loc. 5: 158};
  \node[font=\scriptsize,black] at (517,0.43) {Loc. 50: 152};
\end{tikzpicture}
\caption{Existencia de ProductInventory agregada por las tres ubicaciones
observadas para el producto 317 (593 unidades), junto con sus umbrales de
reorden y seguridad.}
\end{figure}

\section{Fuentes operacionales y selección}

\textbf{Resumen.} AdventureWorks OLTP es una base relacional que registra
pedidos de venta, ofertas, productos, inventario, órdenes de compra y
proveedores. Las tablas de Sales, Production y Purchasing aportan los datos
operacionales para analizar D2 y D7.

\textbf{Modelo conceptual de las fuentes (MERX).} Las relaciones muestran
entidades y atributos que alimentan este subconjunto. Cada fila es una ruta
independiente para evitar cruces visuales.

\small
\begin{longtable}{p{0.20\textwidth}p{0.30\textwidth}p{0.40\textwidth}}
\toprule
\textbf{Relación} & \textbf{Cardinalidad} & \textbf{Atributos relevantes}\\
\midrule
\endhead
Customer -- SalesOrderHeader & Un cliente puede tener muchos pedidos (1:N). &
CustomerID; PersonID/StoreID; SalesOrderID; OrderDate; TerritoryID.\\
SalesOrderHeader -- SalesOrderDetail & Un pedido contiene muchas líneas (1:N). &
SalesOrderID; SalesOrderDetailID; ProductID; SpecialOfferID; OrderQty;
UnitPrice; UnitPriceDiscount; LineTotal.\\
SpecialOffer -- SpecialOfferProduct -- Product & La relación oferta-producto
es N:M, resuelta por entidad asociativa. & SpecialOfferID; ProductID;
DiscountPct; fechas de vigencia.\\
SalesOrderHeader -- CurrencyRate & Una tasa puede referenciarse desde muchos
pedidos; la FK puede ser nula (1:N opcional). & CurrencyRateID; códigos From/To;
AverageRate.\\
SalesOrderHeader -- SalesTerritory & Un territorio puede corresponder a
muchos pedidos; la FK puede ser nula (1:N opcional). & TerritoryID; nombre;
país; grupo/región.\\
ProductSubcategory -- ProductCategory -- Product & Una categoría agrupa
subcategorías y estas productos (1:N). & IDs y nombres; MakeFlag;
FinishedGoodsFlag; ReorderPoint; SafetyStockLevel.\\
Product -- ProductCostHistory & Un producto puede tener varios costos
estándar con períodos de vigencia (1:N). & ProductID; StandardCost; StartDate;
EndDate.\\
Product -- ProductInventory -- Location & El saldo se registra por producto,
ubicación y bin. & ProductID; LocationID; Shelf; Bin; Quantity.\\
Vendor -- PurchaseOrderHeader -- PurchaseOrderDetail -- Product & Un proveedor
tiene órdenes con líneas de productos (1:N); cada línea identifica producto. &
VendorID; PurchaseOrderID; OrderDate; Status; ProductID; OrderQty; UnitPrice;
ReceivedQty; RejectedQty; StockedQty.\\
\bottomrule
\end{longtable}
\normalsize

\textbf{Modelo lógico de las fuentes.} Las claves primarias y foráneas
preservan estas rutas relacionales durante la extracción y la integración.

\small
\begin{longtable}{p{0.16\textwidth}p{0.78\textwidth}}
\toprule
\textbf{Proceso} & \textbf{Entidades, claves y relaciones}\\
\midrule
\endhead
Ventas (D2) &
SalesOrderHeader (PK SalesOrderID; FK CustomerID, TerritoryID,
CurrencyRateID) se relaciona 1:N con SalesOrderDetail (PK SalesOrderID +
SalesOrderDetailID; FK ProductID y SpecialOfferID + ProductID).
Customer, SalesTerritory y CurrencyRate se enlazan mediante sus claves.
SpecialOffer y Product se relacionan mediante SpecialOfferProduct.
ProductCostHistory (PK ProductID + StartDate; FK ProductID) aporta StandardCost
según la fecha de cada venta.\\
Inventario (D7) &
Product (PK ProductID) se relaciona 1:N con ProductInventory (PK ProductID +
LocationID + Shelf + Bin), que se relaciona con Location mediante LocationID.
ProductSubcategory enlaza Product con ProductCategory.\\
Compras (D7) &
PurchaseOrderHeader (PK PurchaseOrderID; FK VendorID) se relaciona 1:N con
PurchaseOrderDetail (PK PurchaseOrderID + PurchaseOrderDetailID; FK ProductID).
Vendor y Product se enlazan mediante VendorID y ProductID.\\
\bottomrule
\end{longtable}
\normalsize

\textbf{Catálogo selectivo.} El catálogo identifica los atributos requeridos
para los análisis y sus claves. PK = clave primaria; FK = clave foránea;
NULL = atributo opcional.

\small
\begin{longtable}{p{0.24\textwidth}p{0.36\textwidth}p{0.15\textwidth}p{0.18\textwidth}}
\toprule
\textbf{Tabla y propósito} & \textbf{Atributos seleccionados (tipo; restricción)} &
\textbf{Uso} & \textbf{Otros atributos}\\
\midrule
\endhead
Sales.\allowbreak SalesOrderHeader: cabecera de venta. &
SalesOrderID (int; PK): pedido; OrderDate (datetime): fecha; Status (tinyint):
estado; CustomerID (int; FK): cliente; TerritoryID (int; FK, NULL): territorio
de venta; OnlineOrderFlag (bit): canal; CurrencyRateID (int; FK, NULL): tasa
asociada. & Fecha, estado realizado, cliente, territorio, canal y moneda del
pedido (D2). & Direcciones, vendedor, impuestos y envío.\\
Sales.\allowbreak SalesOrderDetail: línea de venta. &
SalesOrderID (int; PK/FK): pedido; SalesOrderDetailID (int; PK): línea;
ProductID (int; FK): producto; SpecialOfferID (int; FK compuesta): oferta;
OrderQty (smallint): cantidad; UnitPrice (money): precio unitario;
UnitPriceDiscount (money): descuento aplicado; LineTotal (money; calculado):
importe neto. & Volumen, precio, descuento y venta neta por línea (D2). &
Tracking, rowguid.\\
Sales.\allowbreak SpecialOffer: términos de oferta. &
SpecialOfferID (int; PK): oferta; Description (nvarchar(255)): nombre;
DiscountPct (smallmoney): descuento nominal; Type, Category (nvarchar(50)):
clasificación; StartDate, EndDate (datetime): vigencia; MinQty, MaxQty (int;
MaxQty NULL): límites de cantidad. & Definición y vigencia de oferta (D2). &
rowguid, ModifiedDate.\\
Sales.\allowbreak SpecialOfferProduct: asociación de producto y oferta. &
SpecialOfferID (int; PK/FK): oferta; ProductID (int; PK/FK): producto
elegible. & Verificar la asociación válida producto-oferta (D2). &
rowguid, ModifiedDate.\\
Sales.\allowbreak CurrencyRate: tasas asociadas a pedidos. &
CurrencyRateID (int; PK): tasa; CurrencyRateDate (datetime): fecha de tasa;
FromCurrencyCode, ToCurrencyCode (nchar(3); FK): monedas origen/destino;
AverageRate, EndOfDayRate (money): factores de conversión. & Identificar
moneda y conversión (D2). & Aplicar AverageRate según la regla de conversión
configurada.\\
Sales.\allowbreak Customer: referencia del cliente. &
CustomerID (int; PK): cliente; PersonID (int; FK, NULL): persona asociada;
StoreID (int; FK, NULL): tienda asociada. & Distinguir persona/revendedor y
contar clientes (D2). & Atributos personales de Person/Store.\\
Sales.\allowbreak SalesTerritory: territorio de venta. &
TerritoryID (int; PK): territorio; Name (nvarchar(50)): nombre;
CountryRegionCode (nvarchar(3)): país; Group (nvarchar(50)): región de
negocio. & Filtro geográfico de ventas (D2). & Atributos operacionales.\\
Production.\allowbreak Product: catálogo operacional. &
ProductID (int; PK): producto; Name (nvarchar(50)): nombre;
ProductNumber (nvarchar(25)): código; MakeFlag, FinishedGoodsFlag (bit):
fabricación/terminado; ProductSubcategoryID (int; FK, NULL): subcategoría;
ReorderPoint, SafetyStockLevel (smallint): umbrales de inventario. &
Identificación, selección y umbrales del producto (D2/D7). & Peso, talla,
estilo y especificaciones técnicas.\\
Production.\allowbreak ProductCostHistory: historial de costo estándar. &
ProductID (int; PK/FK): producto; StartDate (datetime; PK): inicio de
vigencia; EndDate (datetime, NULL): fin de vigencia; StandardCost (money):
costo estándar. & Costo aplicable a la fecha del pedido para margen estimado
(D2). & ModifiedDate.\\
Production.\allowbreak ProductSubcategory: clasificación del producto. &
ProductSubcategoryID (int; PK): subcategoría; Name (nvarchar(50)): nombre;
ProductCategoryID (int; FK): categoría superior. & Jerarquía y selección de
componentes (D2/D7). & ModifiedDate.\\
Production.\allowbreak ProductCategory: categoría superior. &
ProductCategoryID (int; PK): categoría; Name (nvarchar(50)): nombre. &
Agrupar productos (D2/D7). & ModifiedDate.\\
Production.\allowbreak ProductInventory: saldo por bin. &
ProductID (int; PK/FK): producto; LocationID (smallint; PK/FK): ubicación;
Shelf (nvarchar(10); PK): estante; Bin (tinyint; PK): contenedor;
Quantity (smallint): unidades disponibles; ModifiedDate (datetime): última
modificación. & Existencia actual por producto/ubicación y captura diaria
DWH (D7). & rowguid.\\
Production.\allowbreak Location: ubicación de inventario/producción. &
LocationID (smallint; PK): ubicación; Name (nvarchar(50)): nombre. &
Desglosar saldo por ubicación (D7). & CostRate, Availability.\\
Purchasing.\allowbreak PurchaseOrderHeader: cabecera de compra. &
PurchaseOrderID (int; PK): orden; OrderDate (datetime): fecha de orden;
Status (smallint): estado; VendorID (int; FK): proveedor. & Período, estado y
proveedor de compra (D7). & Fecha de envío, empleado, impuestos y otros datos
de envío.\\
Purchasing.\allowbreak PurchaseOrderDetail: línea de compra. &
PurchaseOrderID, PurchaseOrderDetailID (int; PK): orden/línea; DueDate
(datetime): fecha prevista; OrderQty (smallint): cantidad pedida; ProductID
(int; FK): material; UnitPrice (money): precio unitario; LineTotal
(numeric(19,4); calculado): importe de línea; ReceivedQty, RejectedQty
(decimal(8,2)): recibidas/rechazadas; StockedQty (decimal(9,2)): ingresadas
al almacén; ModifiedDate (datetime): última modificación. & Precio y
cantidades de compra (D7). & rowguid.\\
Purchasing.\allowbreak Vendor: proveedor. &
BusinessEntityID (int; PK/FK): proveedor; Name (nvarchar(50)): nombre
operativo. & Identificar proveedor en compras (D7). & Cuenta, calificación,
preferencia, contactos y direcciones.\\
\bottomrule
\end{longtable}
\normalsize

\section{Diseño de arquitectura y Data Warehouse}

\textbf{Resumen.} La arquitectura integra las fuentes operacionales,
estandariza datos de ventas, compras e inventario y publica modelos
dimensionales para consulta. El diseño combina almacenamiento por capas,
transformación validada, orquestación y acceso analítico.

\textbf{Diseño del almacén empresarial.} La capa empresarial conserva la
trazabilidad de origen y organiza los dominios integrados con convenciones
comunes. Las claves de origen se mantienen junto a claves sustitutas; fechas,
estados y monedas conservan los valores de negocio para análisis coherente.

\small
\begin{longtable}{p{0.15\textwidth}p{0.35\textwidth}p{0.40\textwidth}}
\toprule
\textbf{Capa} & \textbf{Contenido} & \textbf{Convenciones de integración}\\
\midrule
\endhead
Bronze & Copia de tablas Sales, Production y Purchasing. &
Conservar columnas originales, sistema de origen, clave natural, fecha de
extracción e identificador de lote.\\
Silver & Entidades y procesos validados de ventas, productos, ofertas,
clientes, territorios, compras, proveedores e inventario. &
Aplicar tipos y dominios comunes, preservar estado original y moneda
transaccional, y registrar rechazos con su causa.\\
Gold & Hechos y dimensiones analíticos definidos en esta especificación. &
Usar claves sustitutas con claves de origen para trazabilidad, dimensiones
compartidas y reglas de moneda explícitas.\\
\bottomrule
\end{longtable}
\normalsize

\textbf{Grano y dimensiones.} Se propone un esquema estrella en la capa Gold.
Cada tabla de hechos conserva un proceso y una granularidad únicos. Las
dimensiones de fecha y producto son compartidas. Las claves sustitutas son
enteros; se conserva la clave natural OLTP para trazabilidad.

\small
\begin{longtable}{p{0.20\textwidth}p{0.29\textwidth}p{0.26\textwidth}p{0.17\textwidth}}
\toprule
\textbf{Tabla de hechos / grano} & \textbf{Dimensiones} &
\textbf{Medidas principales} & \textbf{Decisión}\\
\midrule
\endhead
FactSalesLine: una fila por SalesOrderID + SalesOrderDetailID. &
DimDate, DimProduct, DimOffer, DimCustomer, DimTerritory. &
OrderQty, UnitPrice, UnitPriceDiscount, GrossAmount, DiscountAmount,
NetAmount, NetAmountUSD, StandardCostAtSale y GrossProfitUSD. & D2\\
FactPurchaseLine: una fila por PurchaseOrderID + PurchaseOrderDetailID. &
DimDate (roles OrderDate/DueDate), DimProduct, DimVendor. &
OrderQty, UnitPrice, LineTotal, ReceivedQty, RejectedQty y StockedQty. & D7\\
FactInventorySnapshot: una fila por fecha de captura + producto + ubicación;
bins agregados. & DimDate, DimProduct, DimLocation. &
QuantityOnHand, ReorderPointAtSnapshot y SafetyStockLevelAtSnapshot. & D7\\
\bottomrule
\end{longtable}
\normalsize

\textbf{Contenido de dimensiones.} DimDate contiene fecha, día, semana, mes,
trimestre y año. DimProduct conserva ID, nombre/número, categoría,
subcategoría, flags de fabricación/terminado y umbrales. DimOffer contiene
ID, descripción, tipo/categoría, descuento
nominal y vigencia. DimCustomer contiene ID y tipo persona/revendedor.
DimTerritory contiene territorio, país y grupo.
DimVendor contiene clave y nombre operativo. DimLocation contiene clave y
nombre. Las jerarquías analíticas son año--trimestre--mes--día y
categoría--subcategoría--producto; la ubicación se analiza por sus claves y
nombres operacionales.

\textbf{Modelo dimensional resumido.} Las dimensiones se conectan por clave
foránea al hecho indicado; las relaciones son 1:N desde dimensión a hecho.

\begin{center}
\small
\begin{tabular}{p{0.19\textwidth}p{0.24\textwidth}p{0.42\textwidth}}
\toprule
\textbf{Hecho} & \textbf{Proceso} & \textbf{Claves foráneas}\\
\midrule
FactSalesLine & D2: venta por artículo & DateKey, ProductKey, OfferKey,
CustomerKey, TerritoryKey\\
FactPurchaseLine & D7: pedido de material & OrderDateKey, DueDateKey,
ProductKey, VendorKey\\
FactInventorySnapshot & D7: saldo capturado & SnapshotDateKey, ProductKey,
LocationKey\\
\bottomrule
\end{tabular}
\end{center}

\textbf{Transformaciones y carga.}
\begin{enumerate}
  \item Extraer en modo lectura las columnas seleccionadas. Usar extracción
  incremental de pedidos/líneas con claves naturales y marca de modificación;
  reprocesar cambios para evitar duplicados. Capturar ProductInventory
  diariamente de forma completa y consistente.
  \item Conservar una copia de extracción con fecha de lote y recuentos.
  Validar claves, cantidades, descuentos, asociaciones producto-oferta y
  tasas de moneda. Registrar cada fila rechazada con su causa y mantener
  visibles los recuentos de validación.
  \item Generar claves sustitutas y resolver dimensiones antes de los hechos.
  Normalizar fechas al calendario y preservar fecha de negocio. Conservar
  moneda e importe originales, tasa aplicada y estado de conversión a USD.
  \item Calcular en ventas:
  \(\operatorname{GrossAmount}=\operatorname{OrderQty}\cdot\operatorname{UnitPrice}\),
  \(\operatorname{DiscountAmount}=\operatorname{GrossAmount}\cdot\operatorname{UnitPriceDiscount}\) y
  \(\operatorname{NetAmount}=\operatorname{LineTotal}\). Obtener StandardCost
  de ProductCostHistory según la fecha de pedido, convertir el neto a USD y
  calcular \(\operatorname{GrossProfitUSD}=\operatorname{NetAmountUSD}-
  \operatorname{OrderQty}\cdot\operatorname{StandardCost}\). Conciliar el
  neto con LineTotal. Para compras conservar cantidades fuente separadas;
  para inventario sumar bins por producto/ubicación en cada captura diaria.
  \item Publicar dimensiones y hechos en una transacción por lote. Registrar
  fecha de corte, máximas fechas fuente, filas cargadas/rechazadas y
  diferencias de conciliación. Generar alertas de ejecución y completitud
  para cada lote y captura diaria.
\end{enumerate}

\textbf{Historia y reglas de interpretación.} FactInventorySnapshot construye
historia diaria a partir de la primera captura de ProductInventory. Los
umbrales de reorden y seguridad se calculan sobre la existencia total por
producto y se presentan junto con el desglose por ubicación. Para D2, las
ventas netas y los descuentos caracterizan el resultado comercial; el análisis
de margen vincula el costo estándar vigente al período de venta. En compras,
DueDate conserva la fecha prevista y ReceivedQty describe las unidades
recibidas.

\textbf{Componentes de arquitectura.}
\small
\begin{longtable}{p{0.20\textwidth}p{0.71\textwidth}}
\toprule
\textbf{Componente} & \textbf{Responsabilidad}\\
\midrule
\endhead
SQL Server & Proporcionar las tablas operacionales de AdventureWorks.\\
MinIO & Conservar extracciones de origen y datos de las capas del Lake.\\
Spark y dbt & Validar, estandarizar y transformar datos hacia los modelos
analíticos.\\
Dagster & Programar cargas, administrar dependencias y registrar resultados
de ejecución.\\
PostgreSQL & Servir las tablas Gold de hechos y dimensiones.\\
Trino & Consultar las estructuras publicadas para análisis y reportes.\\
Usuarios analíticos & Explorar resultados mediante las visualizaciones
asociadas a D2 y D7.\\
\bottomrule
\end{longtable}
\normalsize

\begin{figure}[htbp]
\centering
\begin{tikzpicture}[node distance=0.45cm,scale=0.9,transform shape]
  \node[box] (source) {SQL Server\\AdventureWorks};
  \node[box,right=of source] (bronze) {MinIO\\extracción Bronze};
  \node[box,right=of bronze] (curated) {Spark / dbt\\validación Silver/Gold};
  \node[box,right=of curated] (gold) {PostgreSQL\\hechos y dimensiones};
  \node[box,right=of gold] (query) {Trino\\consulta};
  \node[box,right=of query] (users) {Usuarios\\analíticos};
  \draw[flow] (source) -- (bronze);
  \draw[flow] (bronze) -- (curated);
  \draw[flow] (curated) -- (gold);
  \draw[flow] (gold) -- (query);
  \draw[flow] (query) -- (users);
\end{tikzpicture}
\caption{Flujo lógico de datos y acceso analítico. Dagster orquesta las cargas
entre componentes.}
\end{figure}

\section{Fuentes de datos y parámetros de configuración}

\begin{itemize}
  \item Datos de ofertas:
  \href{https://raw.githubusercontent.com/microsoft/sql-server-samples/master/samples/databases/adventure-works/oltp-install-script/SpecialOffer.csv}{\texttt{SpecialOffer.csv}}.
  \item Datos de ventas:
  \href{https://raw.githubusercontent.com/microsoft/sql-server-samples/master/samples/databases/adventure-works/oltp-install-script/SalesOrderDetail.csv}{\texttt{SalesOrderDetail.csv}}.
  \item Datos de productos e inventario:
  \href{https://raw.githubusercontent.com/microsoft/sql-server-samples/master/samples/databases/adventure-works/oltp-install-script/Product.csv}{\texttt{Product.csv}} y
  \href{https://raw.githubusercontent.com/microsoft/sql-server-samples/master/samples/databases/adventure-works/oltp-install-script/ProductInventory.csv}{\texttt{ProductInventory.csv}}.
  \item Historial de costos:
  \href{https://raw.githubusercontent.com/microsoft/sql-server-samples/master/samples/databases/adventure-works/oltp-install-script/ProductCostHistory.csv}{\texttt{ProductCostHistory.csv}}.
  \item Datos de compras:
  \href{https://raw.githubusercontent.com/microsoft/sql-server-samples/master/samples/databases/adventure-works/oltp-install-script/PurchaseOrderDetail.csv}{\texttt{PurchaseOrderDetail.csv}}.
\end{itemize}

\textbf{Parámetros de análisis.} La configuración del modelo define moneda
común y tasa de conversión para reportes consolidados; producto y oferta para
el análisis D2; ventana de comparación para ventas; y categoría de materiales,
umbral de inventario y ventana de compras para D7. Los tipos, claves y estados
se corresponden con la estructura operativa de AdventureWorks.

\end{document}
